\section{Related Work}
\label{sec:related_work}

\paragraph{Prompting for Open-Domain QA.}

LLMs can learn various tasks by simply using a few examples as prompts~\cite{originalgpt3}. They've also been shown to answer complex questions by producing step-by-step reasoning (chain-of-thoughts, or CoT) when prompted with a few or zero demonstrations~\cite{cot, zerocot}. Prompting has been applied to open-domain QA~\cite{internet-augmented-qa,recitationlm,Yu2022GenerateRT} but its value in improving retrieval and QA for multi-step open-domain questions remains relatively underexplored.

Recently three approaches have been proposed for multi-step open-domain QA. SelfAsk~\cite{selfask} prompts LLMs to decompose a question into subquestions and answers subquestions by a call to Google Search API. DecomP~\cite{decomp} is a general framework that decomposes a task and delegates sub-tasks to appropriate sub-models. They also decompose questions but delegate retrieval to a BM25-based retriever. Both of these approaches are not developed for CoT reasoning, do not focus on the retrieval problem, and require a single-hop QA model to answer the decomposed questions. Recently proposed ReAct~\cite{react} system frames the problem as generating a sequence of reasoning and action steps. These steps are much more complex, rely on much larger models (PaLM-540B), and require fine-tuning to outperform CoT for multi-step ODQA. Furthermore, none of these works have been shown to be effective for smaller models without any training. While a direct comparison with these approaches is not straightforward (difference in knowledge corpus, LLMs, examples), we find that our ODQA performance is much higher than all their reported numbers where available (\S\ref{sec:exp-results}).


\paragraph{Supervised Multi-Step Open-Domain QA.}
Prior work has explored iterative retrieval for open-domain QA in a fully supervised setting. \citet{multi-step-retriever-reader} proposes an iterative retrieval model that retrieves using a neural query representation and then updates it based on a reading comprehension model's output. \citet{multihop-retriever-odqa} apply similar neural query reformulation idea for multihop open-domain QA. \citet{mpr} extends the widely-used Dense Passage Retrieval (DPR)~\cite{dpr} to multihop setting, which has since been improved by \citet{baleen}. \citet{wikipatretriever} leverages the graph structure induced by the entity links present in Wikipedia paragraphs to perform iterative multi-step retrieval. GoldEn (Gold Entity) retriever~\cite{golden} iteratively generates text queries based on paragraphs retrieved from an off-the-shelf retriever but requires training data for this next query generator. \citet{webgpt} used GPT3 to answer long-form questions by interacting with the browser but relied on human annotations of these interactions. All of these methods rely on supervised training on a large-scale dataset and can not be easily extended to a few-shot setting.
